%!TEX root = ../thesis.tex
\conclusions
    У кваліфікаційній роботі розв’язано актуальну науково-прикладну задачу розробки обчислювально ефективної системи інтерактивної колоризації зображень із можливістю багатоваріантної генерації. За результатами проведених досліджень сформовано такі висновки:

    \begin{enumerate}
        \item У ході роботи було \textbf{повністю виконано завдання щодо аналізу} науково-технічної літератури.
            Дослідження існуючих методів комп'ютерного зору на базі CNN, GAN, Transformer та Diffusion показало, що традиційні моделі зводять задачу до детермінованого відображення, а багатоваріантні аналоги є занадто ресурсоємними для роботи в реальному часі.
            Це обґрунтувало необхідність проєктування гібридної архітектури з використанням SSM.

        \item Завдання з \textbf{підготовки даних та розробки алгоритму генерації підказок виконано в повному обсязі}.
            Сформовано комплексний набір даних для навчання та тестування на базі COCO 2017, Natural Color Dataset та Oxford-102 Flowers.
            Розроблено та програмно реалізовано алгоритм стохастичної генерації колірних підказок на основі градієнтної значущості та семантичних центрів, що дозволило імітувати дії користувача та забезпечити багатоваріантність генерації без ускладнення архітектури мережі.

        \item Успішно \textbf{спроєктовано та реалізовано гібридну нейромережеву архітектуру} B-Mamba, що поєднує локальну точність U-Net із глобальним рецептивним полем двонаправленої Mamba.
            Впровадження спільних ваг дозволило суттєво оптимізувати обчислювальні витрати.
            Для проведення порівняльного аналізу було успішно інтегровано базові моделі інших архітектурних класів, а саме ECCV16, SIGGRAPH17, Pix2Pix, DDColor.

        \item Практичне завдання з \textbf{розробки програмного забезпечення виконано успішно}.
            Створено інтерактивний вебзастосунок із графічним інтерфейсом, який дозволяє користувачу в режимі реального часу керувати процесом колоризації за допомогою підказок.

        \item Завдання щодо \textbf{навчання та комплексного оцінювання розробленої моделі виконано}.
            За результатами об'єктивного тестування модель B-Mamba продемонструвала абсолютну перевагу над аналогами.
            Показник якості $\mathrm{PSNR}$ склав $22.53$ дБ, що на $2.42$ дБ перевищує результат автоматичної моделі ECCV16 та на $1.76$ дБ -- інтерактивної SIGGRAPH17.
            Архітектура забезпечила найкращу перцептивну якість із показниками $\mathrm{LPIPS}$ на рівні $0.143$ та $\mathrm{FID}$ $3.078$ (порівняно з $5.031$ у важкої моделі DDColor).
            Система довела бездоганну точність дотримання користувацьких умов із рекордно низькою похибкою підказок $\mathrm{MSE}_{\mathrm{HINT}}$ на рівні $6.1 \times 10^{-5}$.
            Підтверджено високу апаратну ефективність: час інференсу склав $15.08$ мс при споживанні $187.29$ МБ відеопам'яті, що в $2.3$ раза швидше та майже в $6$ разів економніше порівняно з еталонною моделлю DDColor.

        \item Проведено \textbf{користувацьке дослідження для суб'єктивної оцінки} якості генерації.
            Підтверджено, що запропонована гібридна архітектура ефективно компенсує недоліки суто автоматичної колоризації. 
            Модель отримала високі оцінки загальної реалістичності ($3.509$ з $5$ балів), перевершила базові підходи у здатності точно слідувати просторовим підказкам (перевага у $60.2\%$ випадків) та довела здатність генерувати серію правдоподібних багатоваріантних рішень для складних сцен.
    \end{enumerate}

    \textbf{Перспективи подальших досліджень} у даній предметній області включають:
    \begin{itemize}[label=$\bullet$]
        \item теоретичне дослідження перспективних архітектур, зокрема аналіз можливостей інтеграції дифузійних процесів із SSM (Diffusion-Mamba) для підвищення якості генерації;
        \item удосконалення конвеєра багатоваріантності шляхом розробки допоміжної моделі для автоматичного передбачення природного кольору об'єктів;
        \item розширення набору базових моделей для проведення більш вичерпного та репрезентативного бенчмаркінгу;
        \item адаптацію розробленої архітектури для колоризації відеопотоку із забезпеченням часової узгодженості кадрів;
        \item розширення функціоналу програмного комплексу через впровадження додаткових інструментів для підвищення зручності користувацької взаємодії.
    \end{itemize}

    % У дипломній роботі вирішено актуальну науково-практичну задачу розробки інтерактивної системи колоризації зображень та відео, здатної генерувати фотореалістичні та варіативні результати. У ході виконання роботи були отримані наступні результати:

    % 1.  \textbf{Проведено глибокий аналіз сучасного стану проблеми колоризації.}
    % Встановлено, що класичні методи оптимізації та ранні нейромережеві підходи (на базі CNN з функцією втрат L2) страждають від проблеми "усереднення" кольорів, що призводить до ненасичених, сірих зображень. Порівняльний аналіз показав, що найбільш перспективним напрямком є використання імовірнісних дифузійних моделей (Diffusion Models), які, на відміну від GAN, забезпечують вищу стабільність навчання та кращу деталізацію текстур, хоча і потребують більших обчислювальних ресурсів.

    % 2.  \textbf{Розроблено алгоритм та програмну реалізацію системи інтерактивної колоризації.}
    % За основу взято архітектуру умовної дифузійної моделі з механізмом уваги (Attention). Ключовою особливістю розробленої системи є реалізація механізму багатоваріантної генерації, що дозволяє отримувати декілька правдоподібних варіантів забарвлення для одного вхідного зображення шляхом стохастичного семплювання. Також впроваджено гібридну систему керування процесом, яка приймає як глобальні текстові описи, так і локальні точкові підказки користувача.

    % 3.  \textbf{Виконано експериментальне дослідження ефективності системи.}
    % Тестування на наборі даних COCO-Stuff підтвердило перевагу запропонованого методу над існуючими аналогами (зокрема, методами на основі GAN).
    % \begin{itemize}
    %     \item Значення метрики FID (Fréchet Inception Distance) вдалося знизити до \textbf{24.2}, що на \textbf{25\%} краще за показники популярного методу DeOldify (32.5), що свідчить про вищу реалістичність згенерованих зображень.
    %     \item Суб'єктивна оцінка якості користувачами (Mean Opinion Score) склала \textbf{4.2} бала з 5 можливих, що значно перевищує оцінку повністю автоматичних методів (3.6).
    %     \item Доведено ефективність інтерактивного режиму: додавання всього 2-х точкових підказок дозволяє підвищити точність відновлення кольору (метрика PSNR) у складних сценах в середньому на \textbf{3.3~дБ}.
    % \end{itemize}

    % 4.  \textbf{Визначено межі застосування розробленої системи.}
    % Встановлено, що через специфіку ітеративного процесу дифузії середній час обробки одного зображення складає 4–5 секунд. Це робить систему ідеальною для задач професійної реставрації архівних матеріалів та художньої обробки фотографій, де пріоритетом є якість та контроль, але обмежує її використання в системах реального часу.

    % \textbf{Напрямки подальших досліджень.}
    % Основним вектором розвитку є оптимізація швидкодії алгоритму. Перспективним вбачається перехід до латентних дифузійних моделей (Latent Diffusion Models), що дозволить виконувати генерацію у стиснутому просторі ознак та значно пришвидшити процес. Також для покращення колоризації відео доцільно інтегрувати модулі часової уваги (Temporal Attention) або механізми оптичного потоку безпосередньо в архітектуру нейронної мережі для усунення ефекту мерехтіння без необхідності пост-обробки.
    
% % створюємо Висновки до всієї роботи
% \conclusions
    % Загальні висновки до роботи повинні підсумовувати усі ваші досягнення у 
    % даному напрямку досліджень.

    % За кожним пунктом завдань, поставлених у вступі, у висновках повинен 
    % міститись звіт про виконання: виконано, не виконано, виконано частково (І 
    % чому саме так). Наприклад, якщо першим поставленим завданням у вас іде 
    % <<огляд літератури за тематикою досліджень>>, то на початку висновків ви 
    % повинні зазначити, що <<у ході даної роботи був проведений аналіз 
    % опублікованих джерел за тематикою (...), який показав, що (...)>>. Окрім 
    % простої констатації про виконання ви повинні навести, які саме результати 
    % ви одержали та проінтерпретувати їх з точки зору поставленої задачі, мети 
    % та загальної проблематики.

    % В ідеалі загальні висновки повинні збиратись з висновків до кожного 
    % розділу, але ідеал недосяжний. :) Однак висновки не повинні містити 
    % формул, таблиць та рисунків. Дозволяється (та навіть вітається) 
    % використовувати числа (на кшталт <<розроблена методика дозволяє підвищити 
    % ефективність пустопорожньої балаканини на $2.71\%$>>).

    % Наприкінці висновків необхідно зазначити напрямки подальших досліджень: 
    % куди саме, як вам вважається, необхідно прямувати наступним дослідникам у 
    % даній тематиці.
